\documentclass[12pt]{article}
\usepackage[utf8]{inputenc}
\usepackage[T1]{fontenc}
\usepackage{hyperref}
\title{Kaurna Language Resources for K-12 Curriculum Development}
\author{Compiled by Claude Sonnet 5.5 (claude-sonnet-5-5) for Andr\'as Kornai}
\date{}
\begin{document}
\maketitle

\section{Introduction}
Kaurna (ISO 639-3 zku; Glottocode kaur1267) is a Pama-Nyungan language of the Adelaide Plains in South Australia. It was dormant but is being revived. The Kaurna Warra Pintyanthi (KWP) team at the University of Adelaide produces learning resources. According to a 2023 university announcement, 21 South Australian schools had a registered Kaurna language program. Kaurna is therefore one of the better-resourced revitalization cases.

\section{Primary Educational Resources}
\begin{itemize}
\item University of Adelaide news item on the new Kaurna language website. It describes a learner's guide, dictionary, alphabet book, games, and spelling and pronunciation guidance, and notes the school programs. \url{https://www.adelaide.edu.au/newsroom/news/list/2023/03/09/new-website-helps-grow-kaurna-language}
\item Cosmos Magazine, ``Kaurna online: new website to help people learn the language''. It mentions educational videos in Kaurna, such as a puppet show, animal names, and counting. \url{https://cosmosmagazine.com/people/kaurna-warra-website}
\item InDaily, ``Kaurna online: New website to help learn the language''. \url{https://www.indailysa.com.au/news/archive/2023/03/14/kaurna-online-new-website-to-help-learn-the-language}
\end{itemize}

\section{Linguistic and Reference Materials}
\begin{itemize}
\item Adelaide Digital Library record of a Kaurna language item. \url{https://digital.library.adelaide.edu.au/dspace/handle/2440/93054}
\item Adelaide Digital Library item record. \url{https://digital.library.adelaide.edu.au/items/7100f8d6-deb6-4383-af28-652cac609136/full}
\item An article in the WINHEC journal related to Kaurna. \url{https://journals.uvic.ca/index.php/winhec/article/view/20275/9069}
\end{itemize}
The search results also mention a published dictionary, \emph{Kaurna Warrapiipa}, from Wakefield Press, but I did not see a URL for it.

\section{Religious and Community Resources}
\begin{itemize}
\item Kaurna Warra Pintyanthi (KWP) Team page, University of Adelaide. It describes a committee formalised in 2002 and a team set up in 2012. \url{https://able.adelaide.edu.au/humanities/kaurna-warra-pintyanthi-kwp-team}
\end{itemize}
The search results also mention Kaurna Warra Karrpanthi, an Aboriginal corporation that handles requests for Kaurna names and translations, and Tauondi Aboriginal College, which offers Certificates II and III in Kaurna. I did not see URLs for these.

\section{Digital and Audio Resources}
\begin{itemize}
\item Cosmos Magazine, ``Stories, songs and records: how to bring back a language''. \url{https://cosmosmagazine.com/people/culture/stories-songs-and-records-how-to-bring-back-a-language}
\end{itemize}
The KWP team also runs a YouTube channel and a ``Kaurna for Kids'' Instagram account, per the search results. I did not see their URLs.

\section{Curriculum Development Context}
Kaurna has an active community and university partnership, existing school programs, and online materials. Curriculum developers should work through the KWP team and Kaurna Elders. Note that the reference description lists Kaurna as not known to be taught in schools, but the 2023 university announcement reports 21 registered school programs.
\end{document}